\cotefile{80}
\scanpage{39}
% slanted
\margin{\unsure{correspondent|corresp.|correspond.} $\longrightarrow$\\ à \deleted{\unsure{$\mathfrak S_4\times\{\pm1\}$}}\\ $\unsure{\widetilde{\mathfrak S}_{4}|\widetilde{\mathscr S}_{4}|\widetilde{S}_{4}|\mathfrak S_4}\longrightarrow\mathfrak S_4\times\{\pm1\}\longrightarrow\widetilde{\mathfrak S}_4$}%
où les flèches marquées $\approx$ sont des équivalences,
les flèches marquées $\mathrel{\tikz[baseline=-.45ex]{\draw[->] (0,0)--(.84,0); \draw[dotted] (.08,.02) to[bend left=100] (.12,.15);}}$ sont fidèles, la
flèche $\mathrel{\tikz[baseline=-.45ex]{\draw[->] (0,0)--(.84,0); \draw[dotted] (.81,.02) to[bend right=105] (.66,.17);}}$ est \deleted{\unsure{de}}surjective sur les ens. d'homomorphismes,
la flèche \deleted{verticale} $\mathrel{\tikz[baseline=-.45ex]{\draw[->] (0,.21)--(0,-.21); \node[left,inner sep=.7pt] at (0,.06) {\scriptsize 1};}}$ est \deleted{de}
identité de (oct. g.). La face supérieure du
cube est formée des objets ``orientés'', les quatre
flèches constituant ce carré sont des équivalences ;
la face inférieure est \deleted{\unsure{formée}}\inserted{formée} des objets
non orientés, \deleted{\unsure{dont les sont}} les deux flèches
$\textrm{oct}\longrightarrow\textrm{gr. oct.}\longrightarrow\textrm{gr. oct g.}$ sont des équivalences,
les deux autres sont resp. surjectives sur \inserted{les} hom,
et fidèles. Les quatre flèches de passage
de la théorie orientée vers la non orientée sont
fidèles. La face de gauche est formée des
objets ``ordinaires'' \deleted{ou} ``sphériques'', celle de
droite est formée des objets ``gauches'' ou
``projectifs'', les \inserted{quatre} flèches de passage de l'une
à l'autre sont des équivalences, à l'exception
de $\textrm{oct}\longrightarrow\textrm{oct. gauche}$ (surjectif sur les
hom). La face postérieure du cube est
formée des configurations deux-dimensionnelles,
\deleted{celle} la face antérieure est formée des
configurations $1$-dimensionnelles associées, les
quatre flèches de passage sont des équivalences
sauf $\textrm{oct. g}\longrightarrow\textrm{gr. oct. gauche}$, qui est fidèle
et fait l'objet de l'énoncé (16) plus haut.
% archive sheet number: 38
